% frontmatter.tex -- Module 7 (The Nature of Light): cover, How to use this guide,
% Contents, Module overview.
% Formula list: the formula cards of ch1.tex-ch4.tex (same forms and same badges).
% Experiment list: the Key experiment boxes of ch1.tex-ch4.tex.
\makecover
\IfFileExists{../common/howtouse.tex}{% howtouse.tex -- shared "How to use this guide" pages (same in all four books).
% Input it from frontmatter.tex:   % howtouse.tex -- shared "How to use this guide" pages (same in all four books).
% Input it from frontmatter.tex:   % howtouse.tex -- shared "How to use this guide" pages (same in all four books).
% Input it from frontmatter.tex:   \input{../common/howtouse}
% Numbers in the mini worked example: R = 12 / 0.50 = 24 ohm (checked in Python).
\unnumberedchapter[\faInfoCircle]{How to use this guide}

This guide explains Module~\modulenumber\ (\moduletitle) from the very beginning.
It uses simple English, short sentences and many examples. Every calculation is
shown step by step.

The chapters follow the inquiry questions of the NSW Physics Stage 6 Syllabus
(2017). At the back of the book you will find a formula summary, a glossary
(with Chinese translations), a practice set, full answers and a last-minute
summary. In the PDF, you can click any entry in the contents, any
cross-reference and any question number.

\section*{What the coloured boxes mean}

Each kind of box has its own colour and icon. A colour always means the same
thing, in every book of this series.

\begin{sidebyside}[0.5]
  \begin{bigidea}[before skip=0pt, after skip=0pt, equal height group=legend1, fontupper=\normalsize]
    The main point of a topic in two or three sentences. Read these first.
  \end{bigidea}
\nextside
  \begin{keywords}[before skip=0pt, after skip=0pt, equal height group=legend1]
    \kw{Term}{what the word means, in plain words.}
    New words are also printed in \term{bold} the first time they appear.
  \end{keywords}
\end{sidebyside}

\begin{sidebyside}[0.5]
  \begin{mistakes}[before skip=0pt, after skip=0pt, equal height group=legend2]
    Errors that students often make, and how to avoid them:
    \xmark\ wrong \enspace \cmark\ right.
  \end{mistakes}
\nextside
  \begin{examtip}[before skip=0pt, after skip=0pt, equal height group=legend2]
    How the HSC asks about this topic, and what you must write to get full
    marks.
  \end{examtip}
\end{sidebyside}

\begin{sidebyside}[0.5]
  \begin{toolkit}[before skip=0pt, after skip=0pt, equal height group=legend3]{}
    A short maths skill that you need for the next part, for example
    rearranging a formula.
  \end{toolkit}
\nextside
  \begin{refresher}[before skip=0pt, after skip=0pt, equal height group=legend3]{}
    A short reminder of a Year~11 idea that you need, for example Newton's
    laws.
  \end{refresher}
\end{sidebyside}

\begin{sidebyside}[0.5]
  \begin{tcolorbox}[hsc minor=hscqc, unbreakable, before skip=0pt, after skip=0pt,
      equal height group=legend4, title={\hscboxtitle{\faQuestionCircle}{Quick check}{}}]
    Short questions to test yourself, numbered like \textcolor{hscqc}{\bfseries Q2.3}
    (Chapter~2, question~3). The answers at the back use the same numbers.
  \end{tcolorbox}
\nextside
  \begin{experiment}[unbreakable, before skip=0pt, after skip=0pt, equal height group=legend4]{}
    A practical or historical experiment: aim, apparatus, method, results,
    conclusion and why it matters.
  \end{experiment}
\end{sidebyside}

\begin{formulacard}[unbreakable, before skip=6pt]{Name of the formula}
  \cardformula{V = IR}
  Every formula has a card: the formula, the meaning and SI unit of each
  symbol, when to use it, and a common trap. The badge under the formula tells
  you about the official HSC formula sheet:
  \begin{itemize}
    \item \hscbadgeonsheet\ \ it is printed on the sheet, so you are given it
          in the exam;
    \item \hscbadgenotonsheet\ \ you must learn it, or know how to get it;
    \item \hscbadgecombine{two formulas}\ \ it is two formulas from the sheet
          joined together.
  \end{itemize}
\end{formulacard}

\begin{workedexample}[nonumber, unbreakable]{How every worked example is set out}
  A \qty{12}{V} battery drives a current of \qty{0.50}{A} through a resistor.
  Find the resistance. (Cover everything below the \emph{Solution} line with a
  sheet of paper and try it first.)
  \given{$V = \qty{12}{V}$, \ $I = \qty{0.50}{A}$}
  \find{the resistance $R$}
  \formula{$V = IR$, \ so $R = \dfrac{V}{I}$}
  \substitute{$R = \dfrac{\qty{12}{V}}{\qty{0.50}{A}}$ \ (numbers with their units)}
  \answer{\final{$R = \qty{24}{\ohm}$} \ (with a unit and sensible significant
    figures)}
  \checkstep{is the answer sensible? Units: $\unit{V} \div \unit{A} = \unit{\ohm}$ \cmark}
\end{workedexample}

Each chapter starts with its \textbf{inquiry question} from the syllabus, a
list of \textbf{goals} and the \textbf{syllabus dot points} it covers. It ends
with a short \textbf{summary}.

\section*{How to study with this guide in the last two weeks}

The HSC Physics exam is on \textbf{Thursday 5 November 2026}. You have four
books (Modules 5 to 8). Here is a simple plan.

\begin{table}
  \centering
  \begin{tabular}{@{}L{30mm}L{\dimexpr\linewidth-34mm\relax}@{}}
    \toprule
    \colhead{When} & \colhead{What to do}\\
    \midrule
    Days 1 to 10 & Work through the four books, about two or three days for
      each book. For every chapter, follow steps 1 to 5 below.\\
    Days 11 to 13 & Do the Practice set of each book like a real exam: time
      yourself and use only the formula sheet. Mark your answers with the
      marking criteria. Redo every question you got wrong.\\
    The night before & Read the Last-minute summary of each book and look at
      the formula sheet once more. Then go to bed early.\\
    \bottomrule
  \end{tabular}
\end{table}

For each chapter:
\begin{enumerate}
  \item Read the \textbf{big idea} and the \textbf{key words} (about 10 minutes).
  \item Read the explanation and the \textbf{formula cards}. Say each formula
        ``in words''.
  \item \textbf{Cover the solution} of each worked example and try it yourself.
        Then compare your work with the book, step by step.
  \item Do the \textbf{Quick check} questions. Mark them with the answers at the
        back.
  \item Write every mistake you made in a ``mistakes list''. Read this list
        again before the exam.
\end{enumerate}

\section*{How to use the official formula sheet in the exam}

In the exam you get the NESA data sheet and formulae sheet. It gives many
formulas and constants, but not all of them.

\begin{enumerate}
  \item \textbf{Before the exam}, learn where each formula is on the sheet. The
        sheet has four parts: motion, forces and gravity; waves and
        thermodynamics; electricity and magnetism; quantum, special relativity
        and nuclear. Then you can find a formula quickly.
  \item \textbf{Learn the formulas that are not on the sheet.} Their formula
        cards have the red badge and show you how to derive them.
  \item \textbf{Check what each symbol means} in the question. The same letter
        can mean different things: for example, $T$ is a period in one formula
        and a temperature in another.
  \item \textbf{Change every quantity to SI units first}: km to m, nm to m,
        \si{km.h^{-1}} to \si{m.s^{-1}}, eV to J.
  \item \textbf{Use the values on the data sheet}, for example
        $g = \qty{9.8}{m.s^{-2}}$ and $c = \qty{3.00e8}{m.s^{-1}}$.
  \item \textbf{Show your working} in the same order as the worked examples:
        formula, substitution, answer with a unit.
\end{enumerate}

% Numbers in the mini worked example: R = 12 / 0.50 = 24 ohm (checked in Python).
\unnumberedchapter[\faInfoCircle]{How to use this guide}

This guide explains Module~\modulenumber\ (\moduletitle) from the very beginning.
It uses simple English, short sentences and many examples. Every calculation is
shown step by step.

The chapters follow the inquiry questions of the NSW Physics Stage 6 Syllabus
(2017). At the back of the book you will find a formula summary, a glossary
(with Chinese translations), a practice set, full answers and a last-minute
summary. In the PDF, you can click any entry in the contents, any
cross-reference and any question number.

\section*{What the coloured boxes mean}

Each kind of box has its own colour and icon. A colour always means the same
thing, in every book of this series.

\begin{sidebyside}[0.5]
  \begin{bigidea}[before skip=0pt, after skip=0pt, equal height group=legend1, fontupper=\normalsize]
    The main point of a topic in two or three sentences. Read these first.
  \end{bigidea}
\nextside
  \begin{keywords}[before skip=0pt, after skip=0pt, equal height group=legend1]
    \kw{Term}{what the word means, in plain words.}
    New words are also printed in \term{bold} the first time they appear.
  \end{keywords}
\end{sidebyside}

\begin{sidebyside}[0.5]
  \begin{mistakes}[before skip=0pt, after skip=0pt, equal height group=legend2]
    Errors that students often make, and how to avoid them:
    \xmark\ wrong \enspace \cmark\ right.
  \end{mistakes}
\nextside
  \begin{examtip}[before skip=0pt, after skip=0pt, equal height group=legend2]
    How the HSC asks about this topic, and what you must write to get full
    marks.
  \end{examtip}
\end{sidebyside}

\begin{sidebyside}[0.5]
  \begin{toolkit}[before skip=0pt, after skip=0pt, equal height group=legend3]{}
    A short maths skill that you need for the next part, for example
    rearranging a formula.
  \end{toolkit}
\nextside
  \begin{refresher}[before skip=0pt, after skip=0pt, equal height group=legend3]{}
    A short reminder of a Year~11 idea that you need, for example Newton's
    laws.
  \end{refresher}
\end{sidebyside}

\begin{sidebyside}[0.5]
  \begin{tcolorbox}[hsc minor=hscqc, unbreakable, before skip=0pt, after skip=0pt,
      equal height group=legend4, title={\hscboxtitle{\faQuestionCircle}{Quick check}{}}]
    Short questions to test yourself, numbered like \textcolor{hscqc}{\bfseries Q2.3}
    (Chapter~2, question~3). The answers at the back use the same numbers.
  \end{tcolorbox}
\nextside
  \begin{experiment}[unbreakable, before skip=0pt, after skip=0pt, equal height group=legend4]{}
    A practical or historical experiment: aim, apparatus, method, results,
    conclusion and why it matters.
  \end{experiment}
\end{sidebyside}

\begin{formulacard}[unbreakable, before skip=6pt]{Name of the formula}
  \cardformula{V = IR}
  Every formula has a card: the formula, the meaning and SI unit of each
  symbol, when to use it, and a common trap. The badge under the formula tells
  you about the official HSC formula sheet:
  \begin{itemize}
    \item \hscbadgeonsheet\ \ it is printed on the sheet, so you are given it
          in the exam;
    \item \hscbadgenotonsheet\ \ you must learn it, or know how to get it;
    \item \hscbadgecombine{two formulas}\ \ it is two formulas from the sheet
          joined together.
  \end{itemize}
\end{formulacard}

\begin{workedexample}[nonumber, unbreakable]{How every worked example is set out}
  A \qty{12}{V} battery drives a current of \qty{0.50}{A} through a resistor.
  Find the resistance. (Cover everything below the \emph{Solution} line with a
  sheet of paper and try it first.)
  \given{$V = \qty{12}{V}$, \ $I = \qty{0.50}{A}$}
  \find{the resistance $R$}
  \formula{$V = IR$, \ so $R = \dfrac{V}{I}$}
  \substitute{$R = \dfrac{\qty{12}{V}}{\qty{0.50}{A}}$ \ (numbers with their units)}
  \answer{\final{$R = \qty{24}{\ohm}$} \ (with a unit and sensible significant
    figures)}
  \checkstep{is the answer sensible? Units: $\unit{V} \div \unit{A} = \unit{\ohm}$ \cmark}
\end{workedexample}

Each chapter starts with its \textbf{inquiry question} from the syllabus, a
list of \textbf{goals} and the \textbf{syllabus dot points} it covers. It ends
with a short \textbf{summary}.

\section*{How to study with this guide in the last two weeks}

The HSC Physics exam is on \textbf{Thursday 5 November 2026}. You have four
books (Modules 5 to 8). Here is a simple plan.

\begin{table}
  \centering
  \begin{tabular}{@{}L{30mm}L{\dimexpr\linewidth-34mm\relax}@{}}
    \toprule
    \colhead{When} & \colhead{What to do}\\
    \midrule
    Days 1 to 10 & Work through the four books, about two or three days for
      each book. For every chapter, follow steps 1 to 5 below.\\
    Days 11 to 13 & Do the Practice set of each book like a real exam: time
      yourself and use only the formula sheet. Mark your answers with the
      marking criteria. Redo every question you got wrong.\\
    The night before & Read the Last-minute summary of each book and look at
      the formula sheet once more. Then go to bed early.\\
    \bottomrule
  \end{tabular}
\end{table}

For each chapter:
\begin{enumerate}
  \item Read the \textbf{big idea} and the \textbf{key words} (about 10 minutes).
  \item Read the explanation and the \textbf{formula cards}. Say each formula
        ``in words''.
  \item \textbf{Cover the solution} of each worked example and try it yourself.
        Then compare your work with the book, step by step.
  \item Do the \textbf{Quick check} questions. Mark them with the answers at the
        back.
  \item Write every mistake you made in a ``mistakes list''. Read this list
        again before the exam.
\end{enumerate}

\section*{How to use the official formula sheet in the exam}

In the exam you get the NESA data sheet and formulae sheet. It gives many
formulas and constants, but not all of them.

\begin{enumerate}
  \item \textbf{Before the exam}, learn where each formula is on the sheet. The
        sheet has four parts: motion, forces and gravity; waves and
        thermodynamics; electricity and magnetism; quantum, special relativity
        and nuclear. Then you can find a formula quickly.
  \item \textbf{Learn the formulas that are not on the sheet.} Their formula
        cards have the red badge and show you how to derive them.
  \item \textbf{Check what each symbol means} in the question. The same letter
        can mean different things: for example, $T$ is a period in one formula
        and a temperature in another.
  \item \textbf{Change every quantity to SI units first}: km to m, nm to m,
        \si{km.h^{-1}} to \si{m.s^{-1}}, eV to J.
  \item \textbf{Use the values on the data sheet}, for example
        $g = \qty{9.8}{m.s^{-2}}$ and $c = \qty{3.00e8}{m.s^{-1}}$.
  \item \textbf{Show your working} in the same order as the worked examples:
        formula, substitution, answer with a unit.
\end{enumerate}

% Numbers in the mini worked example: R = 12 / 0.50 = 24 ohm (checked in Python).
\unnumberedchapter[\faInfoCircle]{How to use this guide}

This guide explains Module~\modulenumber\ (\moduletitle) from the very beginning.
It uses simple English, short sentences and many examples. Every calculation is
shown step by step.

The chapters follow the inquiry questions of the NSW Physics Stage 6 Syllabus
(2017). At the back of the book you will find a formula summary, a glossary
(with Chinese translations), a practice set, full answers and a last-minute
summary. In the PDF, you can click any entry in the contents, any
cross-reference and any question number.

\section*{What the coloured boxes mean}

Each kind of box has its own colour and icon. A colour always means the same
thing, in every book of this series.

\begin{sidebyside}[0.5]
  \begin{bigidea}[before skip=0pt, after skip=0pt, equal height group=legend1, fontupper=\normalsize]
    The main point of a topic in two or three sentences. Read these first.
  \end{bigidea}
\nextside
  \begin{keywords}[before skip=0pt, after skip=0pt, equal height group=legend1]
    \kw{Term}{what the word means, in plain words.}
    New words are also printed in \term{bold} the first time they appear.
  \end{keywords}
\end{sidebyside}

\begin{sidebyside}[0.5]
  \begin{mistakes}[before skip=0pt, after skip=0pt, equal height group=legend2]
    Errors that students often make, and how to avoid them:
    \xmark\ wrong \enspace \cmark\ right.
  \end{mistakes}
\nextside
  \begin{examtip}[before skip=0pt, after skip=0pt, equal height group=legend2]
    How the HSC asks about this topic, and what you must write to get full
    marks.
  \end{examtip}
\end{sidebyside}

\begin{sidebyside}[0.5]
  \begin{toolkit}[before skip=0pt, after skip=0pt, equal height group=legend3]{}
    A short maths skill that you need for the next part, for example
    rearranging a formula.
  \end{toolkit}
\nextside
  \begin{refresher}[before skip=0pt, after skip=0pt, equal height group=legend3]{}
    A short reminder of a Year~11 idea that you need, for example Newton's
    laws.
  \end{refresher}
\end{sidebyside}

\begin{sidebyside}[0.5]
  \begin{tcolorbox}[hsc minor=hscqc, unbreakable, before skip=0pt, after skip=0pt,
      equal height group=legend4, title={\hscboxtitle{\faQuestionCircle}{Quick check}{}}]
    Short questions to test yourself, numbered like \textcolor{hscqc}{\bfseries Q2.3}
    (Chapter~2, question~3). The answers at the back use the same numbers.
  \end{tcolorbox}
\nextside
  \begin{experiment}[unbreakable, before skip=0pt, after skip=0pt, equal height group=legend4]{}
    A practical or historical experiment: aim, apparatus, method, results,
    conclusion and why it matters.
  \end{experiment}
\end{sidebyside}

\begin{formulacard}[unbreakable, before skip=6pt]{Name of the formula}
  \cardformula{V = IR}
  Every formula has a card: the formula, the meaning and SI unit of each
  symbol, when to use it, and a common trap. The badge under the formula tells
  you about the official HSC formula sheet:
  \begin{itemize}
    \item \hscbadgeonsheet\ \ it is printed on the sheet, so you are given it
          in the exam;
    \item \hscbadgenotonsheet\ \ you must learn it, or know how to get it;
    \item \hscbadgecombine{two formulas}\ \ it is two formulas from the sheet
          joined together.
  \end{itemize}
\end{formulacard}

\begin{workedexample}[nonumber, unbreakable]{How every worked example is set out}
  A \qty{12}{V} battery drives a current of \qty{0.50}{A} through a resistor.
  Find the resistance. (Cover everything below the \emph{Solution} line with a
  sheet of paper and try it first.)
  \given{$V = \qty{12}{V}$, \ $I = \qty{0.50}{A}$}
  \find{the resistance $R$}
  \formula{$V = IR$, \ so $R = \dfrac{V}{I}$}
  \substitute{$R = \dfrac{\qty{12}{V}}{\qty{0.50}{A}}$ \ (numbers with their units)}
  \answer{\final{$R = \qty{24}{\ohm}$} \ (with a unit and sensible significant
    figures)}
  \checkstep{is the answer sensible? Units: $\unit{V} \div \unit{A} = \unit{\ohm}$ \cmark}
\end{workedexample}

Each chapter starts with its \textbf{inquiry question} from the syllabus, a
list of \textbf{goals} and the \textbf{syllabus dot points} it covers. It ends
with a short \textbf{summary}.

\section*{How to study with this guide in the last two weeks}

The HSC Physics exam is on \textbf{Thursday 5 November 2026}. You have four
books (Modules 5 to 8). Here is a simple plan.

\begin{table}
  \centering
  \begin{tabular}{@{}L{30mm}L{\dimexpr\linewidth-34mm\relax}@{}}
    \toprule
    \colhead{When} & \colhead{What to do}\\
    \midrule
    Days 1 to 10 & Work through the four books, about two or three days for
      each book. For every chapter, follow steps 1 to 5 below.\\
    Days 11 to 13 & Do the Practice set of each book like a real exam: time
      yourself and use only the formula sheet. Mark your answers with the
      marking criteria. Redo every question you got wrong.\\
    The night before & Read the Last-minute summary of each book and look at
      the formula sheet once more. Then go to bed early.\\
    \bottomrule
  \end{tabular}
\end{table}

For each chapter:
\begin{enumerate}
  \item Read the \textbf{big idea} and the \textbf{key words} (about 10 minutes).
  \item Read the explanation and the \textbf{formula cards}. Say each formula
        ``in words''.
  \item \textbf{Cover the solution} of each worked example and try it yourself.
        Then compare your work with the book, step by step.
  \item Do the \textbf{Quick check} questions. Mark them with the answers at the
        back.
  \item Write every mistake you made in a ``mistakes list''. Read this list
        again before the exam.
\end{enumerate}

\section*{How to use the official formula sheet in the exam}

In the exam you get the NESA data sheet and formulae sheet. It gives many
formulas and constants, but not all of them.

\begin{enumerate}
  \item \textbf{Before the exam}, learn where each formula is on the sheet. The
        sheet has four parts: motion, forces and gravity; waves and
        thermodynamics; electricity and magnetism; quantum, special relativity
        and nuclear. Then you can find a formula quickly.
  \item \textbf{Learn the formulas that are not on the sheet.} Their formula
        cards have the red badge and show you how to derive them.
  \item \textbf{Check what each symbol means} in the question. The same letter
        can mean different things: for example, $T$ is a period in one formula
        and a temperature in another.
  \item \textbf{Change every quantity to SI units first}: km to m, nm to m,
        \si{km.h^{-1}} to \si{m.s^{-1}}, eV to J.
  \item \textbf{Use the values on the data sheet}, for example
        $g = \qty{9.8}{m.s^{-2}}$ and $c = \qty{3.00e8}{m.s^{-1}}$.
  \item \textbf{Show your working} in the same order as the worked examples:
        formula, substitution, answer with a unit.
\end{enumerate}
}{}
\tableofcontents

\unnumberedchapter[\faMap]{Module overview}

% Small marks for the "on the sheet?" column (same words, icons and colours as \frow).
\newcommand\ovSheetYes{\textcolor{hscformula}{\footnotesize\bfseries\faCheck\ Yes}}
\newcommand\ovSheetNo{\textcolor{hscmistake}{\footnotesize\bfseries\faTimes\ Learn it}}
\newcommand\ovSheetCombine{\textcolor{hsctiptext}{\footnotesize\bfseries\faLink\ Combine}}
% \ovPad{height}{depth}: extra room in a table row for a tall formula
\newcommand\ovPad[2]{\rule[-#2]{0pt}{0pt}\rule{0pt}{#1}}

\section*{What this module is about}

Module~7 asks one big question: \emph{what is light?} Newton (about 1700) said light is a
stream of tiny particles. Huygens said it is a wave. Diffraction, interference and
polarisation later showed that light behaves as a wave.

In the 1860s, Maxwell showed that light is an electromagnetic (EM) wave: a wave of electric
and magnetic fields. Spectra (light spread out into its colours) tell us which elements a
lamp or a star contains.

Around 1900, the glow of hot objects and the photoelectric effect did not fit the wave model.
Planck and Einstein explained them: light energy comes in packets, called photons. So light
acts as a wave \emph{and} as a stream of particles.

Every observer measures the same speed for light in a vacuum. From this, Einstein showed that moving
clocks run slow, moving objects are shorter, and mass is a form of energy ($E = mc^2$).

\section*{The inquiry questions}

Each chapter answers one inquiry question of the syllabus.

\begin{enumerate}[label=\textbf{IQ\arabic*}, leftmargin=3.2em, itemsep=3pt]
  \item \emph{What is light?}\hfill \cref{ch1:chap}\\
        Maxwell and EM waves, measuring the speed of light, and spectra of elements and stars.
  \item \emph{What evidence supports the classical wave model of light, and what
        predictions can be made using this model?}\hfill \cref{ch2:chap}\\
        Newton and Huygens, diffraction, interference, gratings and polarisation.
  \item \emph{What evidence supports the particle model of light, and what does this
        evidence mean for the quantum model of light?}\hfill \cref{ch3:chap}\\
        Black bodies, Planck's quanta, the photoelectric effect and Einstein's photons.
  \item \emph{How does the behaviour of light affect concepts of time, space and
        matter?}\hfill \cref{ch4:chap}\\
        Einstein's postulates, time dilation, length contraction, $p_\mathrm{v}$ and
        $E = mc^2$.
\end{enumerate}

\section*{How the parts fit together}

\Cref{ov:fig:map} shows how the four chapters are linked. Start at the top and follow the
arrows.

\begin{figure}
  \centering
  \begin{tikzpicture}[x=1cm, y=1cm]
    \node[concept=37mm] (A) at (0,0)
      {\textbf{Light is an EM wave}\\Maxwell, 1860s\\[1pt]{\footnotesize\color{hscmuted}Chapter 1}};
    \node[concept=37mm] (D) at (-5.4,-3.6)
      {\textbf{Wave model}\\diffraction, interference, polarisation\\[1pt]
       {\footnotesize\color{hscmuted}Chapter 2}};
    \node[concept=37mm] (C) at (0,-3.6)
      {\textbf{Spectra}\\identify elements; study stars\\[1pt]
       {\footnotesize\color{hscmuted}Chapter 1}};
    \node[concept=37mm] (B) at (5.4,-3.6)
      {\textbf{The speed of light}\\measured; now fixed, and it defines the metre\\[1pt]
       {\footnotesize\color{hscmuted}Chapter 1}};
    \node[concept=37mm] (E) at (-5.4,-7.4)
      {\textbf{Particle model}\\black bodies, the photoelectric effect, \mbox{$E = hf$}\\[1pt]
       {\footnotesize\color{hscmuted}Chapter 3}};
    \node[concept=37mm] (F) at (0,-7.4)
      {\textbf{Wave--particle duality}\\light acts as a wave \emph{and} as particles\\[1pt]
       {\footnotesize\color{hscmuted}Chapters 2 and 3}};
    \node[concept=37mm] (G) at (5.4,-7.4)
      {\textbf{Special relativity}\\time dilation, length contraction, $p_\mathrm{v}$,
       \mbox{$E = mc^2$}\\[1pt]{\footnotesize\color{hscmuted}Chapter 4}};
    \draw[link] (A.west) -| node[linklabel, pos=0.3, above] {a transverse wave} (D.north);
    \draw[link] (A.south) -- node[linklabel] {a spectrum of wavelengths} (C.north);
    \draw[link] (A.east) -| node[linklabel, pos=0.3, above] {predicts the speed $c$} (B.north);
    \draw[link] (D.south) -- node[linklabel, align=center] {fails for hot objects\\and the
      photoelectric effect} (E.north);
    \draw[link] (D.south east) -- node[linklabel, sloped] {waves} (F.north west);
    \draw[link] (E.east) -- node[linklabel, above=1pt] {photons} (F.west);
    \draw[link] (B.south) -- node[linklabel, align=center] {the same for\\every observer} (G.north);
  \end{tikzpicture}
  \caption{How the ideas in Module 7 are linked.}\label{ov:fig:map}
\end{figure}

\section*{Formulas at a glance}

\Cref{ov:tab:formulas} lists every formula card in this book. The last column is about the
official HSC formula sheet: \ovSheetYes\ = printed on the sheet; \ovSheetNo\ = learn it or
derive it; \ovSheetCombine\ = join two formulas from the sheet. The \emph{Formula sheet
explained} chapter (page~\pageref{fs:start}) gives more detail.

{\small
\renewcommand\arraystretch{1.3}\arrayrulecolor{hscrule}%
\setlength\LTpre{4pt}\setlength\LTpost{2pt}%
\begin{longtable}{@{}C{44mm}@{\hspace{2mm}}L{72mm}@{\hspace{2mm}}C{8mm}@{\hspace{2mm}}C{20mm}@{}}
  \caption{Every formula card in Module 7, in chapter order.}\label{ov:tab:formulas}\\
  \colhead{Formula} & \colhead{What it is for} & \colhead{Ch} & \colhead{On the sheet?}\\\hline
  \endfirsthead
  \colhead{Formula} & \colhead{What it is for} & \colhead{Ch} & \colhead{On the sheet?}\\\hline
  \endhead
  \ovPad{3.2ex}{2.6ex}$c = \dfrac{1}{\sqrt{\mu_0\varepsilon_0}}$ & Maxwell's predicted speed of EM waves & 1 & \ovSheetNo\\\hline
  $c = f\lambda$ & frequency and wavelength of an EM wave in a vacuum or air & 1 & \ovSheetCombine\\\hline
  \ovPad{3.0ex}{1.8ex}$c = \dfrac{d}{t}$ & the speed of light from a distance and a time & 1 & \ovSheetNo\\\hline
  \ovPad{3.0ex}{1.8ex}$\lambda_\mathrm{max} = \dfrac{b}{T}$ & Wien's law: peak wavelength and temperature & 1, 3 & \ovSheetYes\\\hline
  \ovPad{3.0ex}{1.8ex}$n_\mathrm{x} = \dfrac{c}{v_\mathrm{x}}$ & the speed of light in water or glass & 2 & \ovSheetYes\\\hline
  $d\sin\theta = m\lambda$ & bright fringes of a double slit or a grating & 2 & \ovSheetYes\\\hline
  $d\sin\theta = \left(m - \tfrac{1}{2}\right)\lambda$ & dark fringes of a double slit ($m = 1, 2, 3, \ldots$) & 2 & \ovSheetNo\\\hline
  $\delta = \left|S_2P - S_1P\right|$ & path difference from the distances to two sources: bright (loud) or dark (quiet)? & 2 & \ovSheetNo\\\hline
  $y = L\tan\theta$ & where a fringe lands on the screen & 2 & \ovSheetNo\\\hline
  \ovPad{3.0ex}{1.8ex}$\Delta y = \dfrac{\lambda L}{d}$ & fringe spacing (double slit, small angles) & 2 & \ovSheetNo\\\hline
  \ovPad{3.0ex}{1.8ex}$d = \dfrac{1}{N}$ & grating spacing from lines per metre & 2 & \ovSheetNo\\\hline
  $I = \tfrac{1}{2}I_0$ & unpolarised light through the first polariser & 2 & \ovSheetNo\\\hline
  $I = I_\mathrm{max}\cos^2\theta$ & Malus' law: polarised light through a polariser & 2 & \ovSheetYes\\\hline
  $E = hf$ & the energy of one quantum (photon) & 3 & \ovSheetYes\\\hline
  \ovPad{3.0ex}{1.8ex}$E = \dfrac{hc}{\lambda}$ & photon energy from the wavelength ($E = hf$ card) & 3 & \ovSheetCombine\\\hline
  $K_\mathrm{max} = qV_0$ & stopping voltage and maximum kinetic energy & 3 & \ovSheetCombine\\\hline
  $K_\mathrm{max} = hf - \phi$ & Einstein's photoelectric equation & 3 & \ovSheetYes\\\hline
  \ovPad{3.0ex}{1.8ex}$\phi = hf_0 \quad \lambda_0 = \dfrac{hc}{\phi}$ & threshold frequency and threshold wavelength & 3 & \ovSheetNo\\\hline
  \ovPad{3.0ex}{1.8ex}$\dfrac{N}{t} = \dfrac{P}{hf} = \dfrac{P\lambda}{hc}$ & photons each second from a lamp or laser & 3 & \ovSheetNo\\\hline
  \ovPad{3.0ex}{1.8ex}$V_0 = \dfrac{h}{q}f - \dfrac{\phi}{q}$ & reading $h$, $f_0$ and $\phi$ from a $V_0$--$f$ graph & 3 & \ovSheetNo\\\hline
  \ovPad{3.4ex}{5.0ex}$t = \dfrac{t_0}{\sqrt{\left(1 - \dfrac{v^2}{c^2}\right)}}$ & time dilation: a moving clock runs slow & 4 & \ovSheetYes\\\hline
  \ovPad{3.6ex}{2.6ex}$l = l_0\sqrt{\left(1 - \dfrac{v^2}{c^2}\right)}$ & length contraction along the motion & 4 & \ovSheetYes\\\hline
  \ovPad{3.4ex}{5.0ex}$p_\mathrm{v} = \dfrac{m_0 v}{\sqrt{\left(1 - \dfrac{v^2}{c^2}\right)}}$ & relativistic momentum at any speed & 4 & \ovSheetYes\\\hline
  $E = mc^2$ & energy from lost mass: fusion, annihilation, burning & 4 & \ovSheetYes\\\hline
  \ovPad{3.0ex}{1.8ex}$\dfrac{\Delta m}{\Delta t} = \dfrac{P}{c^2}$ & mass lost each second by the Sun or a power station & 4 & \ovSheetCombine\\\hline
\end{longtable}}

% Layout (book QA): keep the heading and its intro with the start of the table.
\Needspace{14\baselineskip}
\section*{Key experiments}

These are the experiments in this module (\cref{ov:tab:experiments}). Learn what each one
showed: questions often ask you to recount one, or to assess its validity, reliability or
accuracy. A star ($\ast$) marks an investigation that the syllabus says you must carry out.

{\small
\renewcommand\arraystretch{1.2}\arrayrulecolor{hscrule}%
\setlength\LTpre{4pt}\setlength\LTpost{2pt}%
\begin{longtable}{@{}L{52mm}@{\hspace{3mm}}L{72mm}@{\hspace{3mm}}L{19mm}@{}}
  \caption{Key experiments in Module 7.}\label{ov:tab:experiments}\\
  \colhead{Experiment} & \colhead{What it showed} & \colhead{Where}\\\hline
  \endfirsthead
  \colhead{Experiment} & \colhead{What it showed} & \colhead{Where}\\\hline
  \endhead
  Hertz's radio waves (1887) & EM waves exist and travel at about $c$ &
    \cref{ch1:sec:maxwell}\\\hline
  Galileo's lanterns (about 1638) & light is far too fast to time by hand &
    \cref{ch1:sec:experiments}\\\hline
  Rømer: the eclipses of Io (1676) & light takes time to travel (its speed is finite) &
    \cref{ch1:sec:experiments}\\\hline
  Fizeau's toothed wheel (1849) & the first value of $c$ measured on Earth &
    \cref{ch1:sec:experiments}\\\hline
  Measuring $c$ with a microwave oven $\ast$ & a modern school method: $c = f\lambda$ &
    \cref{ch1:sec:experiments}\\\hline
  Spectra of gases, a filament and sunlight $\ast$ & continuous, emission and absorption
    spectra & \cref{ch1:sec:experiments}\\\hline
  Foucault: light in water (1850) & light is slower in water, so Newton was wrong &
    \cref{ch2:sec:experiments}\\\hline
  Diffraction of laser light $\ast$ & a narrower slit spreads the light more &
    \cref{ch2:sec:experiments}\\\hline
  Young's double slit; gratings $\ast$ & interference: light is a wave; measuring $\lambda$ &
    \cref{ch2:sec:experiments}\\\hline
  Testing Malus' law $\ast$ & $I = I_\mathrm{max}\cos^2\theta$; light is transverse &
    \cref{ch2:sec:experiments}\\\hline
  Black body curves (1899--1900) & classical physics fails; Planck's quanta &
    \cref{ch3:sec:experiments}\\\hline
  The photoelectric effect & a threshold frequency; $h$ and $\phi$ from a graph &
    \cref{ch3:sec:experiments}\\\hline
  Michelson and Morley (1887) & no aether (a null result): supports the postulates &
    \cref{ch4:sec:experiments}\\\hline
  Cosmic-ray muons (1963) & many more reach sea level than expected: time dilation &
    \cref{ch4:sec:experiments}\\\hline
  Atomic clocks on aeroplanes (1971) & flown clocks gain or lose time as predicted &
    \cref{ch4:sec:experiments}\\\hline
  Particle accelerators (1964, 1977) & particles never reach $c$; fast muons live longer &
    \cref{ch4:sec:experiments}\\\hline
  Type Ia supernovae & distant ones fade more slowly: time dilation &
    \cref{ch4:sec:experiments}\\\hline
  Thought experiments (light clock; train) & time dilation; ``at the same time'' depends on
    the observer & \cref{ch4:sec:time}\\\hline
\end{longtable}}
